% SGA 3，第 VII_A 讲， 英文重建稿——第 1 轮。
% 范围：权威底本局部页 37--40, 段落 5.4--5.5.3,
% 止于其前： 节 6 局部页 40.
% 合并阅读本页码： 490--493. 再版印刷页码： 479--482.
% 本范围内原版页边定位： 452--454.
% 权威底本： Polo--Gille Exp7A-13oct24.pdf,
% SHA-256 D46DB8C79D0385643B6EED4B96900613248375F9B73B4202478A4F57FD3E8522.
% 曾参考 Jacob Reinhold 的 CC BY 4.0 Markdown 作为比对谱系。

\medskip
\noindent\textbf{5.4.}
\phantomsection
\label{sga3:VIIA:5.4}

% 权威底本： Exp7A 局部页 37; 合并本页 490; 印刷页 479.
\(U_p(\mathfrak g)\) 的泛性质表明，\(U_p(\mathfrak g)\) 关于 \(\mathfrak g\) 具有函子性。每个受限 \(p\)-Lie 代数同态
\[
  \varphi:\mathfrak g\longrightarrow\mathfrak h
\]
诱导出唯一的含幺代数同态 \(U_p(\varphi)\)，使得
\[
  j_{\mathfrak h}\circ\varphi
  =
  U_p(\varphi)\circ j_{\mathfrak g}.
\]
下面给出一些例子。

\smallskip
\noindent\textit{a)}
若 \(\mathfrak h=0\)，则 \(U_p(\mathfrak h)\) 等同于基环，而 \(U_p(\varphi)\) 是一个代数同态
\[
  \varepsilon_{\mathfrak g}:
  U_p(\mathfrak g)\longrightarrow R
\]
称为\emph{增广}。

\smallskip
\noindent\textit{b)}
现在令 \(\mathfrak h=\mathfrak g^\circ\) 为 \(\mathfrak g\) 的反 Lie 代数。于是 \(\mathfrak g^\circ\) 与 \(\mathfrak g\) 具有相同的底层模和符号 \(p\) 次幂，而其括号是 \(\mathfrak g\) 的括号的负值。可以将 \(U_p(\mathfrak g^\circ)\) 等同于 \(U_p(\mathfrak g)\) 的反代数。此外，从 \(\mathfrak g\) 到 \(\mathfrak g^\circ\) 的同构
\[
  x\longmapsto-x
\]
诱导出同构
\[
  c_{\mathfrak g}:
  U_p(\mathfrak g)
  \longrightarrow
  U_p(\mathfrak g^\circ)
  \simeq
  U_p(\mathfrak g)^\circ.
\]
我们称 \(c_{\mathfrak g}\) 为 \(U_p(\mathfrak g)\) 的\emph{对极}。

\smallskip
\noindent\textit{c)}
最后，设 \(\mathfrak f\) 和 \(\mathfrak g\) 是两个受限 \(p\)-Lie 代数，令 \(\mathfrak h=\mathfrak f\times\mathfrak g\) 为它们的积。其底层 \(R\)-模是直积，且
\[
  [(x,y),(x',y')]
  =
  ([x,x'],[y,y']),
  \qquad
  (x,y)^{(p)}
  =
  (x^{(p)},y^{(p)}).
\]
若
\[
  h_1:\mathfrak f\to\mathfrak k,
  \qquad
  h_2:\mathfrak g\to\mathfrak k
\]
是满足下式的受限 \(p\)-Lie 代数同态：
\[
  [h_1(x),h_2(y)]=0
  \quad
  (x\in\mathfrak f,\ y\in\mathfrak g),
\]
则
\[
  h_1+h_2:(x,y)\longmapsto h_1(x)+h_2(y)
\]
是受限 \(p\)-Lie 代数同态。反之，每个同态 \(\mathfrak f\times\mathfrak g\to\mathfrak k\) 都具有这种形式。这通过一个泛性质刻画了 \(\mathfrak f\times\mathfrak g\)。

例如，映射
\[
  h_1:x\longmapsto j_{\mathfrak f}(x)\otimes1,
  \qquad
  h_2:y\longmapsto1\otimes j_{\mathfrak g}(y)
\]
诱导出从 \(\mathfrak f\times\mathfrak g\) 到下述代数的底层受限 \(p\)-Lie 代数的同态 \(h_1+h_2\)：
\[
  U_p(\mathfrak f)\otimes U_p(\mathfrak g).
\]
\(\mathfrak f\times\mathfrak g\) 与受限包络代数的泛性质表明，该同态延拓为同构
\[
  \varphi:
  U_p(\mathfrak f\times\mathfrak g)
  \xrightarrow{\sim}
  U_p(\mathfrak f)\otimes U_p(\mathfrak g).
\]

\medskip
\noindent\textbf{定义。}
\phantomsection
\label{sga3:VIIA:5.4-comultiplication-definition}
若 \(\mathfrak f=\mathfrak g\)，则从 \(\mathfrak g\) 到 \(\mathfrak g\times\mathfrak g\) 的对角映射
\[
  \delta:x\longmapsto(x,x)
\]
诱导出从 \(U_p(\mathfrak g)\) 到 \(U_p(\mathfrak g\times\mathfrak g)\) 的同态。将此同态与 \(\varphi\) 的复合记为
\[
  \Delta_{\mathfrak g}:
  U_p(\mathfrak g)
  \longrightarrow
  U_p(\mathfrak g)\otimes U_p(\mathfrak g)
\]
。
\footnote{因此，对于每个 \(x\in\mathfrak g\)，有 \(\Delta_{\mathfrak g}(x)=x\otimes1+1\otimes x\)；特别地，\(\Delta_{\mathfrak g}\) 是余交换的。}

于是容易看出，\(\Delta_{\mathfrak g}\) 连同 \(U_p(\mathfrak g)\) 的乘法，使 \(U_p(\mathfrak g)\) 成为群型 \(R\)-余代数（参见 \SGARefLink{sga3:VIIA:3.2}{3.2}），其增广为 \(\varepsilon_{\mathfrak g}\)，对极为 \(c_{\mathfrak g}\)。

\medskip
\noindent\textbf{5.5.}
\phantomsection
\label{sga3:VIIA:5.5}
\footnote{编者将原文的 \S5.4.1 重组为此处的 \S5.5。\SGARefLink{sga3:VIIA:5.5.1}{命题 5.5.1} 合并了 \SGARefLink{sga3:VIIA:5}{第~5~节} 和 \SGARefLink{sga3:VIIA:3.2.3}{命题~3.2.3} 的结果，并纳入了原文的 \SGARefLink{sga3:VIIA:7.3}{引理~7.3}；命题 5.5.3(ii) 的证明扩充了原文 \SGARefLink{sga3:VIIA:7.4}{定理~7.4} 中 \((i)\Rightarrow(ii)\) 的证明。}

% 权威底本： Exp7A 局部页 38; 合并本页 491; 印刷页 480.
设 \(S\) 是特征为 \(p\) 的概形。先设 \(\mathcal U\) 是群 \(\mathcal O_S\)-余代数，并令
\[
  G=\operatorname{Spec}^{*}\mathcal U.
\]
我们在 \SGARefLink{sga3:VIIA:3.2.3}{3.2.3} 中看到，对于每个 \(T\to S\)，\((\operatorname{Lie}G)(T)\) 是 \(\Gamma(T,\mathcal U_T)\) 中由本原元素构成的 Lie 子代数。若 \(x\) 是本原的，则
\[
  \Delta(x^p)
  =
  x^p\otimes1+1\otimes x^p,
\]
因为当 \(0<i<p\) 时有 \(\binom pi\equiv0\pmod p\)。因此 \(x^p\) 仍是本原的。由 \SGARefLink{sga3:VIIA:5.1}{5.1} 和 \SGARefLink{sga3:VIIA:5.2}{5.2}，映射 \(x\mapsto x^p\) 使 \((\operatorname{Lie}G)(T)\) 成为 \(\mathcal O(T)\) 上的受限 \(p\)-Lie 代数。

现在设 \(\mathcal L\) 是 \(\mathcal O_S\) 上的拟凝聚受限 \(p\)-Lie 代数。当 \(V\) 遍历 \(S\) 的开子集时，上面在
\[
  U_p\bigl(\Gamma(V,\mathcal L)\bigr)
\]
上定义的群余代数结构，在相伴层 \(\mathcal U_p(\mathcal L)\) 上诱导出群 \(\mathcal O_S\)-余代数结构。此外，对于每个 \(S\)-概形 \(T\)，有同构
\[
  \mathcal U_p(\mathcal L_T)
  \xrightarrow{\sim}
  \mathcal U_p(\mathcal L)_T.
\]

令 \(\operatorname{Prim}\mathcal U_p(\mathcal L)\) 表示如下子预层：它将每个开子集 \(V\) 对应到 \(\mathcal U_p(\mathcal L)(V)\) 中的本原元素；容易看出它是一个层。当 \(V\) 遍历 \(S\) 的开子集时，复合映射
\begin{center}
  \phantomsection
  \label{sga3:VIIA:5.5-primitive-sheaf-sequence}
  % SGA3-VII-DIAG-052
  \(\displaystyle
\Gamma(V,\mathcal L)
  \xrightarrow{j}
\operatorname{Prim}\mathrm U_p\!\bigl(\Gamma(V,\mathcal L)\bigr)
  \longrightarrow
\operatorname{Prim}\mathcal U_p(\mathcal L)(V)
\)

\end{center}
定义了态射
\[
  \mathcal L
  \longrightarrow
  \operatorname{Prim}\mathcal U_p(\mathcal L),
\]
，仍记为 \(j\) 或 \(j_{\mathcal L}\)。它又定义了 \(\mathcal O_S\) 上受限 \(p\)-Lie 代数的态射
\[
  \mathcal W(\mathcal L)
  \longrightarrow
  \operatorname{Prim}
    \mathcal W\bigl(\mathcal U_p(\mathcal L)\bigr)
\]
（参见 \SGARefLink{sga3:VIIA:3.2.3}{3.2.3}）。

\medskip
\noindent\textbf{命题 5.5.1。}
\phantomsection
\label{sga3:VIIA:5.5.1}
\textit{设 \(\mathcal L\) 是 \(\mathcal O_S\) 上的受限 \(p\)-Lie 代数，且作为 \(\mathcal O_S\)-模局部自由。则 \(j_{\mathcal L}\) 诱导出 \(\mathcal O_S\) 上受限 \(p\)-Lie 代数的同构}
\[
  \mathcal W(\mathcal L)
  \xrightarrow{\sim}
  \operatorname{Prim}
    \mathcal W\bigl(\mathcal U_p(\mathcal L)\bigr).
\]

\begin{proof}
设 \(T\) 是 \(S\)-概形。鉴于等同
\[
  \mathcal U_p(\mathcal L_T)
  =
  \mathcal U_p(\mathcal L)_T,
\]
我们必须证明
\[
  \Gamma(T,\mathcal L_T)
  \longrightarrow
  \operatorname{Prim}
    \Gamma\bigl(T,\mathcal U_p(\mathcal L_T)\bigr)
\]
是双射。以 \(T\) 代替 \(S\)，便归结到 \(T=S\)；此时只须证明层态射
\[
  j_{\mathcal L}:
  \mathcal L
  \longrightarrow
  \operatorname{Prim}\mathcal U_p(\mathcal L)
\]
是同构。问题在 \(S\) 上是局部的，因此可设 \(S=\operatorname{Spec}R\)，且 \(\mathcal L\) 是与一个具有全序基 \((x_\alpha)\) 的 \(R\) 上受限 \(p\)-Lie 代数 \(L\) 相伴的层。

如 \SGARefLink{sga3:VIIA:5.3.3}{5.3.3} 中一样，令 \(z_\alpha\) 为 \(x_\alpha\) 在 \(U=U_p(L)\) 中的像。对于满足 \(0\leq n_\alpha\leq p-1\) 且只有有限多个非零项的族 \(n=(n_\alpha)\)，令
\[
  z^{(n)}
  =
  \prod_\alpha\frac{z_\alpha^{n_\alpha}}{n_\alpha!},
\]
，其中乘积按递增顺序取。由于
\[
  \Delta(z_\alpha)
  =
  z_\alpha\otimes1+1\otimes z_\alpha,
\]
得到
\[
  \Delta(z^{(n)})
  =
  \sum_r z^{(n-r)}\otimes z^{(r)},
\]
其中求和遍历满足下式的 \(r\) 所成的有限集：
\[
  0\leq r_\alpha\leq n_\alpha
  \quad\text{对于每个 }\alpha.
\]
\(z^{(n)}\) 和 \(z^{(n)}\otimes z^{(m)}\) 分别构成 \(U\) 和 \(U\otimes U\) 的一组基。由此可知，对于 \(u\in U\)，等式
\[
  \Delta(u)=u\otimes1+1\otimes u
\]
成立当且仅当 \(u\) 是 \(z_\alpha\) 的线性组合。命题得证。
\end{proof}

\medskip
\noindent\textbf{注 5.5.2。}
\phantomsection
\label{sga3:VIIA:5.5.2}
回顾 \SGARefLink{sga3:VIIA:3.2.2}{3.2.2} 和 \SGARefLink{sga3:VIIA:3.2.3}{3.2.3}，群 \(S\)-函子
\[
  G=\operatorname{Spec}^{*}\mathcal U_p(\mathcal L)
\]
是甚好的，并且
\[
  \operatorname{Lie}G
  =
  \operatorname{Prim}
    \mathcal W\bigl(\mathcal U_p(\mathcal L)\bigr).
\]
因此，\SGARefLink{sga3:VIIA:5.5.1}{命题 5.5.1} 表明，\(j_{\mathcal L}\) 诱导出同构
\[
  \mathcal W(\mathcal L)
  \xrightarrow{\sim}
  \operatorname{Lie}G.
\]

若进一步有 \(\mathcal L\) 局部自由且秩有限，则由 \SGARefLink{sga3:VIIA:5.3.3}{5.3.3}，\(\mathcal U_p(\mathcal L)\) 在 \(\mathcal O_S\) 上有限局部自由。因此 \(\operatorname{Spec}^{*}\mathcal U_p(\mathcal L)\) 由 \(S\)-群
\[
  \mathfrak G_p(\mathcal L)
  =
  \operatorname{Spec}\mathcal U_p(\mathcal L)^{*}
\]
表示（参见 \SGARefLink{sga3:VIIA:3.2.2.1}{3.2.2.1}）。我们得到下面更精确的陈述。

\medskip
\noindent\textbf{命题 5.5.3。}
\phantomsection
\label{sga3:VIIA:5.5.3}
\textit{设 \(\mathcal L\) 是 \(\mathcal O_S\) 上的受限 \(p\)-Lie 代数，且作为 \(\mathcal O_S\)-模局部自由、秩有限。令}
\[
  \mathcal A=\mathcal U_p(\mathcal L)^{*},
  \qquad
  G=\mathfrak G_p(\mathcal L)=\operatorname{Spec}\mathcal A.
\]
\textit{则：}

\smallskip
\noindent
\textit{\SGAInlineRefTarget{sga3:VIIA:5.5.3:item:i}(i) \(j_{\mathcal L}\) 诱导出同构}
\[
  \mathcal W(\mathcal L)
  \xrightarrow{\sim}
  \operatorname{Lie}G
\]
\textit{，这是 \(\mathcal O_S\) 上受限 \(p\)-Lie 代数的同构。}

\smallskip
\noindent
\textit{\SGAInlineRefTarget{sga3:VIIA:5.5.3:item:ii}(ii) 令 \(\mathcal I\) 为 \(\mathcal A\) 的增广理想，并令 \(\omega_G=\mathcal I/\mathcal I^2\)（参见 II，\SGARefLink{sga3:ii:remark:4-11-4}{4.11.4}）。则}
\[
  \omega_G
  \simeq
  \mathcal L^{*}
  =
  \operatorname{Hom}_{\mathcal O_S}
    (\mathcal L,\mathcal O_S).
\]
\textit{特别地，\(\omega_G\) 局部自由且秩有限，并且 \(\omega_{G/S}^{*}\simeq\mathcal L\)。}

\begin{proof}
\SGARefLink{sga3:VIIA:5.5.3:item:i}{(i)} 由 \SGARefLink{sga3:VIIA:5.5.2}{注~5.5.2} 推出。下面证明 \SGARefLink{sga3:VIIA:5.5.3:item:ii}{(ii)}。以 \(\eta_{\mathcal U}\) 和 \(\varepsilon_{\mathcal U}\) 表示 \(\mathcal U=\mathcal U_p(\mathcal L)\) 的单位和增广，以 \(\eta_{\mathcal A}\) 和 \(\varepsilon_{\mathcal A}\) 表示 \(\mathcal A\) 的单位和增广。令
\[
  \mathcal J=\ker\varepsilon_{\mathcal U}.
\]
则
\begin{equation}
  \tag{1}\label{sga3:vii:proposition:5-5-3:equation:1}
  \mathcal U
  =
  \eta_{\mathcal U}(\mathcal O_S)\oplus\mathcal J.
\end{equation}
令 \(\delta\) 为由下图定义的态射
\begin{center}
  \phantomsection
  \label{sga3:VIIA:5.5.3-augmentation-kernel-delta-square}
  % SGA3-VII-DIAG-053
  \begin{tikzcd}[row sep=3.2em,column sep=4.2em]
\mathcal J \arrow[r,"\delta"] \arrow[d,swap,"\tau"]
  & \mathcal J\otimes\mathcal J \\
\mathcal U \arrow[r,"\Delta"]
  & \mathcal U\otimes\mathcal U \arrow[u,swap,"\pi"]
\end{tikzcd}

\end{center}
其中 \(\tau\) 和 \(\pi\) 是由分解 \SGARefLink{sga3:vii:proposition:5-5-3:equation:1}{(1)} 得到的包含映射和投影。存在正合列
\begin{center}
  \phantomsection
  \label{sga3:VIIA:5.5.3-primitive-exact-sequence}
  % SGA3-VII-DIAG-054
  \(\displaystyle
(\ast)\qquad
0 \longrightarrow \mathcal L^*
  \xrightarrow{j_{\mathcal L}}
\mathcal J
  \xrightarrow{\delta}
\mathcal J\otimes\mathcal J .
\)

\end{center}

由 \SGARefLink{sga3:VIIA:5.3.3}{5.3.3}，\(\mathcal O_S\)-模 \(\mathcal J/\mathcal L\) 局部自由，而由 \SGARefLink{sga3:VIIA:5.5.1}{5.5.1}，上面的列在每次基变换后仍然正合。因此，\(\delta\) 诱导出 \(\mathcal J/\mathcal L\) 到 \(\mathcal J\otimes\mathcal J\) 的一个局部直和项 \(Q\) 的同构。由对偶性，此正合列遂成为
\begin{center}
  \phantomsection
  \label{sga3:VIIA:5.5.3-dual-exact-sequence}
  % SGA3-VII-DIAG-055
  \(\displaystyle
(\ast\ast)\qquad
0 \longleftarrow \mathcal L
  \xleftarrow{{}^{t}\!j_{\mathcal L}}
\mathcal I
  \xleftarrow{{}^{t}\!\delta}
\mathcal I\otimes\mathcal I .
\)

\end{center}

% 权威底本： Exp7A 局部页 40; 合并本页 493; 印刷页 482.
分解 \SGARefLink{sga3:vii:proposition:5-5-3:equation:1}{(1)} 通过对偶对应于
\begin{equation}
  \tag{2}\label{sga3:vii:proposition:5-5-3:equation:2}
  \mathcal A
  =
  \eta_{\mathcal A}(\mathcal O_S)\oplus\mathcal I.
\end{equation}
\(\Delta\) 的转置是乘法
\[
  m_{\mathcal A}:
  \mathcal A\otimes\mathcal A
  \longrightarrow
  \mathcal A.
\]
由于 \(\mathcal I\) 是 \(\mathcal A\) 的理想，\(m_{\mathcal A}\) 将 \(\mathcal I\otimes\mathcal I\) 映入 \(\mathcal I\)。更精确地，分解 \SGARefLink{sga3:vii:proposition:5-5-3:equation:2}{(2)} 给出交换方图
\begin{center}
  \phantomsection
  \label{sga3:VIIA:5.5.3-augmentation-ideal-multiplication-square}
  % SGA3-VII-DIAG-056
  \begin{tikzcd}[row sep=3.1em,column sep=4.2em]
\mathcal I
  & \mathcal I\otimes\mathcal I \arrow[l,swap,"m'"] \arrow[d,"{}^{t}\!\pi"] \\
\mathcal A \arrow[u,"{}^{t}\!\tau"]
  & \mathcal A\otimes\mathcal A . \arrow[l,swap,"m_{\mathcal A}"]
\end{tikzcd}

\end{center}
这表明 \(m_{\mathcal A}\) 在 \(\mathcal I\otimes\mathcal I\) 上的限制 \(m'\) 是 \(\delta\) 的转置。因此，对偶正合列给出
\[
  \mathcal I/\mathcal I^2
  \simeq
  \mathcal L^{*},
\]
命题得证。
\end{proof}

% 下一游标： Exp7A 局部页 40, 合并阅读本页 493,
% 再版印刷页 482, 位于 节 6.
